% Bagian: counterfactuals
% Bab: minimal-change-semantics
% Subbagian: antecedent-strengthening

\documentclass[../../../include/open-logic-section]{subfiles}

\begin{document}

\olfileid[id]{cnt}{min}{agg}

\olsection{Penguatan Anteseden}

``Memperkuat anteseden'' mengacu pada inferensi $!A \lif !C \Entails (!A
\land !B) \lif !C$. Inferensi ini valid bagi kondisional material, tetapi
tidak valid bagi kontrafaktual. Andaikan benar bahwa seandainya saya menggesek
korek api ini, korek tersebut akan menyala. (Artinya, tidak ada masalah dengan
korek api atau permukaan kotaknya, saya tidak akan mematahkan korek itu, dan
sebagainya.) Namun, tidak benar bahwa seandainya saya menggesek korek api ini
di luar angkasa, korek tersebut akan menyala. Jadi, inferensi berikut tidak
valid:
\begin{quote}
  Seandainya korek api itu digesek, korek tersebut akan menyala.

  Oleh karena itu, seandainya korek api itu digesek di luar angkasa, korek
  tersebut akan menyala.
\end{quote}

Penjelasan Lewis--Stalnaker tentang kondisional menerangkan hal ini: dunia
terdekat tempat saya menggesek korek api dan melakukannya di luar angkasa
jauh lebih berbeda dari dunia aktual daripada dunia terdekat tempat saya
menggesek korek api. Jadi, meskipun benar bahwa korek api menyala di dunia
yang disebut terakhir, hal itu tidak benar di dunia yang disebut pertama.
Dan memang demikianlah semestinya.

\begin{ex}
  Semantik bola membuat inferensi tersebut tidak valid, yakni kita mempunyai
  $p \cif r \Entails/ (p \land q) \cif r$. Perhatikan model $\mModel{M}
  = \tuple{W, O, V}$ dengan $W = \{w, w_1, w_2\}$, $O_w = \{\{w\}, \{w,
  w_1\}, \{w, w_1, w_2\}\}$, $V(p) = \{w_1, w_2\}$, $V(q) = \{w_2\}$,
  dan $V(r) = \{w_1\}$. Terdapat bola yang memuat dunia-$p$, yaitu $S =
  \{w, w_1\}$, dan $p \lif r$ benar pada semua dunia di dalamnya, sehingga
  $\mSat{M}{p \cif r}[w]$. Terdapat pula bola yang memuat dunia-$(p \land q)$,
  yaitu $S' = \{w, w_1, w_2\}$, tetapi $\mSat/{M}{(p \land q) \lif r}[w_2]$,
  sehingga $\mSat/{M}{(p \land q) \cif r}[w]$ (lihat
  \olref{fig:antecedent-strengthening}).

\begin{figure}
\begin{center}
\begin{tikzpicture}[layerwidth=1.5,scale=.6]\tiny
  \spheresystem{3}
  \propositionintersect{-10}{3}{30}{5}
  \begin{scope}
    \clip plot[smooth,tension=1] coordinates
          {(0,4.5) (30:.95) (4.5,-3.5)} -- (4.5,4.5) ;
    \spherefill{2}
  \end{scope}
  \draw[proposition] plot[dashed,smooth,tension=1] coordinates
          {(0,4.5) (30:.95) (4.5,-3.5)} ;
  \proposition{30}{2}{30}{5}
  \path (0,0) node[world] {$w$};
  \spherepos{30}{2}{node[world] {$w_1$}}
  \spherepos{-10}{3}{node[world] {$w_2$}}
  \spherepos{-10}{4}{node {$q$}}
  \spherepos{30}{4}{node {$r$}}
  \draw (1,4.5) node[below] {$p$};
\end{tikzpicture}
\caption{Model tandingan bagi penguatan anteseden}
\ollabel{fig:antecedent-strengthening}
\end{center}
\end{figure}
\end{ex}

\end{document}
